\section{Formal power series: definition and basic properties}

\subsection{Convention}

Fix a commutative ring $K$. (For example, $K$ can be $\mathbb{Z}$ or
$\mathbb{Q}$ or $\mathbb{C}$.)

\subsection{The definition of formal power series}

\begin{definition}
\label{def.fps.fps}
\lean{AlgebraicCombinatorics.FPS.fpsEquivSeq}
\leantarget
\leanok
A \emph{formal power series} (or, short, \emph{FPS}) in $1$
indeterminate over $K$ means a sequence $\left(a_{0},a_{1},a_{2}%
,\ldots\right) = \left(a_{n}\right)_{n\in\mathbb{N}} \in K^{\mathbb{N}}$ of
elements of $K$.
\end{definition}

\begin{definition}
\label{def.fps.ops}
\lean{AlgebraicCombinatorics.FPS.coeff_add_fps, AlgebraicCombinatorics.FPS.coeff_sub_fps, AlgebraicCombinatorics.FPS.coeff_smul_fps, AlgebraicCombinatorics.FPS.coeff_mul_fps, AlgebraicCombinatorics.FPS.coeff_C_fps}
\leantarget
\leanok
\textbf{(a)} The \emph{sum} of two FPSs $\mathbf{a}=\left(
a_{0},a_{1},a_{2},\ldots\right)$ and $\mathbf{b}=\left(b_{0},b_{1}%
,b_{2},\ldots\right)$ is defined to be the FPS%
\[
\left(a_{0}+b_{0},\ \ a_{1}+b_{1},\ \ a_{2}+b_{2},\ \ \ldots\right).
\]
It is denoted by $\mathbf{a}+\mathbf{b}$. \medskip

\textbf{(b)} The \emph{difference} of two FPSs $\mathbf{a}=\left(a_{0}%
,a_{1},a_{2},\ldots\right)$ and $\mathbf{b}=\left(b_{0},b_{1},b_{2}%
,\ldots\right)$ is defined to be the FPS%
\[
\left(a_{0}-b_{0},\ \ a_{1}-b_{1},\ \ a_{2}-b_{2},\ \ \ldots\right).
\]
It is denoted by $\mathbf{a}-\mathbf{b}$. \medskip

\textbf{(c)} If $\lambda\in K$ and if $\mathbf{a}=\left(a_{0},a_{1}%
,a_{2},\ldots\right)$ is an FPS, then we define an FPS
\[
\lambda\mathbf{a}:=\left(\lambda a_{0},\lambda a_{1},\lambda a_{2}%
,\ldots\right).
\]

\textbf{(d)} The \emph{product} of two FPSs $\mathbf{a}=\left(a_{0}%
,a_{1},a_{2},\ldots\right)$ and $\mathbf{b}=\left(b_{0},b_{1},b_{2}%
,\ldots\right)$ is defined to be the FPS $\left(c_{0},c_{1},c_{2}%
,\ldots\right)$, where
\begin{align*}
c_{n} & =\sum_{i=0}^{n}a_{i}b_{n-i}=\sum_{\substack{\left(i,j\right)
\in\mathbb{N}^{2};\\i+j=n}}a_{i}b_{j}\\
& =a_{0}b_{n}+a_{1}b_{n-1}+a_{2}b_{n-2}+\cdots+a_{n}b_{0}%
\ \ \ \ \ \ \ \ \ \ \text{for each }n\in\mathbb{N}.
\end{align*}
This product is denoted by $\mathbf{a}\cdot\mathbf{b}$ or just by
$\mathbf{ab}$. \medskip

\textbf{(e)} For each $a\in K$, we define $\underline{a}$ to be the FPS
$\left(a,0,0,0,\ldots\right)$. An FPS of the form $\underline{a}$ for some
$a\in K$ (that is, an FPS $\left(a_{0},a_{1},a_{2},\ldots\right)$
satisfying $a_{1}=a_{2}=a_{3}=\cdots=0$) is said to be \emph{constant}.
\medskip

\textbf{(f)} The set of all FPSs (in $1$ indeterminate over $K$) is denoted
$K\left[\left[x\right]\right]$.
\end{definition}

\begin{definition}
\label{def.fps.coeff}
\lean{AlgebraicCombinatorics.FPS.fps_coeff_mk}
\leantarget
\leanok
If $n\in\mathbb{N}$, and if $\mathbf{a}=\left(
a_{0},a_{1},a_{2},\ldots\right) \in K\left[\left[x\right]\right]$ is
an FPS, then we define an element $\left[x^{n}\right]\mathbf{a}\in K$ by
\[
\left[x^{n}\right]\mathbf{a}:=a_{n}.
\]
This is called the \emph{coefficient of }$x^{n}$\emph{ in }$\mathbf{a}$, or
the $n$\emph{-th coefficient} of $\mathbf{a}$, or the $x^{n}$%
\emph{-coefficient} of $\mathbf{a}$.
\end{definition}

\begin{theorem}
\label{thm.fps.ring}
\lean{AlgebraicCombinatorics.FPS.smul_mul_fps, AlgebraicCombinatorics.FPS.smul_eq_C_mul_fps}
\leantarget
\leanok
\textbf{(a)} The set $K\left[\left[x\right]\right]$
is a commutative ring (with its operations $+$, $-$ and $\cdot$ defined in
Definition \ref{def.fps.ops}). Its zero and its unity are the FPSs
$\underline{0}=\left(0,0,0,\ldots\right)$ and $\underline{1}=\left(
1,0,0,0,\ldots\right)$. This means, concretely, that the following facts hold:

\begin{enumerate}
\item \emph{Commutativity of addition:} We have $\mathbf{a}+\mathbf{b}%
=\mathbf{b}+\mathbf{a}$ for all $\mathbf{a},\mathbf{b}\in K\left[\left[
x\right]\right]$.

\item \emph{Associativity of addition:} We have $\mathbf{a}+\left(
\mathbf{b}+\mathbf{c}\right) =\left(\mathbf{a}+\mathbf{b}\right)
+\mathbf{c}$ for all $\mathbf{a},\mathbf{b},\mathbf{c}\in K\left[\left[
x\right]\right]$.

\item \emph{Neutrality of zero:} We have $\mathbf{a}+\underline{0}%
=\underline{0}+\mathbf{a}=\mathbf{a}$ for all $\mathbf{a}\in K\left[\left[
x\right]\right]$.

\item \emph{Subtraction undoes addition:} Let $\mathbf{a},\mathbf{b}%
,\mathbf{c}\in K\left[\left[x\right]\right]$. We have $\mathbf{a}%
+\mathbf{b}=\mathbf{c}$ if and only if $\mathbf{a}=\mathbf{c}-\mathbf{b}$.

\item \emph{Commutativity of multiplication:} We have $\mathbf{ab}%
=\mathbf{ba}$ for all $\mathbf{a},\mathbf{b}\in K\left[\left[x\right]
\right]$.

\item \emph{Associativity of multiplication:} We have $\mathbf{a}\left(
\mathbf{bc}\right) =\left(\mathbf{ab}\right)\mathbf{c}$ for all
$\mathbf{a},\mathbf{b},\mathbf{c}\in K\left[\left[x\right]\right]$.

\item \emph{Distributivity:} We have%
\[
\mathbf{a}\left(\mathbf{b}+\mathbf{c}\right) =\mathbf{ab}+\mathbf{ac}%
\ \ \ \ \ \ \ \ \ \ \text{and}\ \ \ \ \ \ \ \ \ \ \left(\mathbf{a}%
+\mathbf{b}\right)\mathbf{c}=\mathbf{ac}+\mathbf{bc}%
\]
for all $\mathbf{a},\mathbf{b},\mathbf{c}\in K\left[\left[x\right]
\right]$.

\item \emph{Neutrality of one:} We have $\mathbf{a}\cdot\underline{1}%
=\underline{1}\cdot\mathbf{a}=\mathbf{a}$ for all $\mathbf{a}\in K\left[
\left[x\right]\right]$.

\item \emph{Annihilation:} We have $\mathbf{a}\cdot\underline{0}%
=\underline{0}\cdot\mathbf{a}=\underline{0}$ for all $\mathbf{a}\in K\left[
\left[x\right]\right]$.
\end{enumerate}

\textbf{(b)} Furthermore, $K\left[\left[x\right]\right]$ is a
$K$-module (with its scaling being the map that sends each $\left(
\lambda,\mathbf{a}\right) \in K\times K\left[\left[x\right]\right]$
to the FPS $\lambda\mathbf{a}$ defined in Definition \ref{def.fps.ops}
\textbf{(c)}). Its zero vector is $\underline{0}$. Concretely, this means that:

\begin{enumerate}
\item[10.] \emph{Associativity of scaling:} We have $\lambda\left(
\mu\mathbf{a}\right)  =\left(  \lambda\mu\right)  \mathbf{a}$ for all
$\lambda,\mu\in K$ and $\mathbf{a}\in K\left[\left[x\right]\right]$.

\item[11.] \emph{Left distributivity:} We have $\lambda\left(  \mathbf{a}%
+\mathbf{b}\right)  =\lambda\mathbf{a}+\lambda\mathbf{b}$ for all $\lambda\in
K$ and $\mathbf{a},\mathbf{b}\in K\left[\left[x\right]\right]$.

\item[12.] \emph{Right distributivity:} We have $\left(  \lambda+\mu\right)
\mathbf{a}=\lambda\mathbf{a}+\mu\mathbf{a}$ for all $\lambda,\mu\in K$ and
$\mathbf{a}\in K\left[\left[x\right]\right]$.

\item[13.] \emph{Neutrality of one:} We have $1\mathbf{a}=\mathbf{a}$ for all
$\mathbf{a}\in K\left[\left[x\right]\right]$.

\item[14.] \emph{Left annihilation:} We have $0\mathbf{a}=\underline{0}$ for
all $\mathbf{a}\in K\left[\left[x\right]\right]$.

\item[15.] \emph{Right annihilation:} We have $\lambda\underline{0}%
=\underline{0}$ for all $\lambda\in K$.
\end{enumerate}

\textbf{(c)} We have $\lambda\left(\mathbf{a}\cdot\mathbf{b}\right)
=\left(\lambda\mathbf{a}\right)\cdot\mathbf{b}=\mathbf{a}\cdot\left(
\lambda\mathbf{b}\right)$ for all $\lambda\in K$ and $\mathbf{a}%
,\mathbf{b}\in K\left[\left[x\right]\right]$. \medskip

\textbf{(d)} Finally, we have $\lambda\mathbf{a}=\underline{\lambda}%
\cdot\mathbf{a}$ for all $\lambda\in K$ and $\mathbf{a}\in K\left[\left[
x\right]\right]$.
\end{theorem}

\begin{proof}
Most parts of Theorem \ref{thm.fps.ring}
are straightforward to verify. Let us check the associativity of multiplication.

Let $\mathbf{a},\mathbf{b},\mathbf{c}\in K\left[\left[x\right]\right]$.
We must prove that $\mathbf{a}\left(\mathbf{bc}\right) = \left(
\mathbf{ab}\right)\mathbf{c}$. Let $n\in\mathbb{N}$. Consider the two
equalities
\begin{align*}
\left[x^{n}\right]\left(\left(\mathbf{ab}\right)\mathbf{c}\right)
& =\sum_{j=0}^{n}\underbrace{\left[x^{n-j}\right]\left(\mathbf{ab}%
\right)}_{=\sum_{i=0}^{n-j}\left[x^{i}\right]\mathbf{a}%
\cdot\left[x^{n-j-i}\right]\mathbf{b}}\cdot\left[x^{j}\right]\mathbf{c}\\
& =\sum_{j=0}^{n}\ \ \sum_{i=0}^{n-j}\left[x^{i}\right]\mathbf{a}%
\cdot\left[x^{n-j-i}\right]\mathbf{b}\cdot\left[x^{j}\right]\mathbf{c}%
\end{align*}
and
\begin{align*}
\left[x^{n}\right]\left(\mathbf{a}\left(\mathbf{bc}\right)\right)
& =\sum_{i=0}^{n}\left[x^{i}\right]\mathbf{a}\cdot\underbrace{\left[
x^{n-i}\right]\left(\mathbf{bc}\right)}_{=\sum_{j=0}%
^{n-i}\left[x^{n-i-j}\right]\mathbf{b}\cdot\left[x^{j}\right]
\mathbf{c}}\\
& =\sum_{i=0}^{n}\ \ \sum_{j=0}^{n-i}\left[x^{i}\right]\mathbf{a}%
\cdot\left[x^{n-i-j}\right]\mathbf{b}\cdot\left[x^{j}\right]
\mathbf{c}.
\end{align*}
The right hand sides are equal since
\[
\sum_{j=0}^{n}\ \ \sum_{i=0}^{n-j}=\sum_{\substack{\left(i,j\right)
\in\mathbb{N}^{2};\\i+j\leq n}}=\sum_{i=0}^{n}\ \ \sum_{j=0}^{n-i}
\]
and $n-j-i=n-i-j$. Thus $\left[x^{n}\right]\left(\left(\mathbf{ab}\right)\mathbf{c}\right)
=\left[x^{n}\right]\left(\mathbf{a}\left(\mathbf{bc}\right)\right)$
for each $n\in\mathbb{N}$, which entails $\left(\mathbf{ab}\right)
\mathbf{c}=\mathbf{a}\left(\mathbf{bc}\right)$.

The remaining claims of Theorem \ref{thm.fps.ring} follow by similar
(and easier) coefficient-wise verifications.
In the Lean formalization, the commutative ring and module structures are
provided by Mathlib's existing instances on power series.
\end{proof}

\begin{lemma}
\label{lem.fps.coeff_mul_range}
\lean{AlgebraicCombinatorics.FPS.coeff_mul_fps'}
\leanhelper
\leanok
The product formula can be written as a sum over a range:
\[
\left[x^{n}\right]\left(\mathbf{ab}\right)
= \sum_{i=0}^{n}\left[x^{i}\right]\mathbf{a}\cdot\left[x^{n-i}\right]\mathbf{b}.
\]
\end{lemma}

\begin{proof}
This rewrites the antidiagonal sum from Definition~\ref{def.fps.ops}
as a range sum using the bijection $(i, n-i) \leftrightarrow i \in \{0, \ldots, n\}$.
\end{proof}

\subsection{Essentially finite sums and summable families}

\begin{definition}
\label{def.infsum.essfin}
\lean{AlgebraicCombinatorics.FPS.EssentiallyFinite}
\leantarget
\leanok
\textbf{(a)} A family $\left(a_{i}\right)_{i\in I}\in K^{I}$
of elements of $K$ is said to be \emph{essentially finite} if all
but finitely many $i\in I$ satisfy $a_{i}=0$ (in other words, if the set
$\left\{i\in I\ \mid\ a_{i}\neq0\right\}$ is finite). \medskip

\textbf{(b)} Let $\left(a_{i}\right)_{i\in I}\in K^{I}$ be an essentially
finite family of elements of $K$. Then, the infinite sum $\sum_{i\in I}a_{i}$
is defined to equal the finite sum $\sum_{\substack{i\in I;\\a_{i}\neq0}}a_{i}$.
Such an infinite sum is said to be \emph{essentially finite}.
\end{definition}

\begin{definition}
\label{def.fps.summable}
\lean{AlgebraicCombinatorics.FPS.SummableFPS, AlgebraicCombinatorics.FPS.summableFPSSum}
\leantarget
\leanok
A (possibly infinite) family $\left(\mathbf{a}_{i}\right)_{i\in I}$
of FPSs is said to be \emph{summable} (or \emph{entrywise essentially finite})
if
\[
\text{for each }n\in\mathbb{N}\text{, all but finitely many }i\in I\text{
satisfy }\left[x^{n}\right]\mathbf{a}_{i}=0.
\]
In this case, the sum $\sum_{i\in I}\mathbf{a}_{i}$ is defined to be the FPS
with
\[
\left[x^{n}\right]\left(\sum_{i\in I}\mathbf{a}_{i}\right)
=\underbrace{\sum_{i\in I}\left[x^{n}\right]\mathbf{a}_{i}}%
_{\substack{\text{an essentially}\\\text{finite sum}}}
\ \ \ \ \ \ \ \ \ \ \text{for all }n\in\mathbb{N}\text{.}
\]
\end{definition}

\begin{proposition}
\label{prop.fps.summable.sub}
\lean{AlgebraicCombinatorics.FPS.summableFPS_subfamily}
\leantarget
\leanok
Let $\left(\mathbf{a}_{i}\right)_{i\in I}$ be
a summable family of FPSs. Then, any subfamily of $\left(\mathbf{a}%
_{i}\right)_{i\in I}$ is summable as well.
\end{proposition}

\begin{proof}
\leanok
Let $J$ be a subset of $I$.
We must prove that the subfamily $\left(\mathbf{a}_{i}\right)_{i\in J}$ is summable.

Let $n\in\mathbb{N}$. Then, all but finitely many $i\in I$ satisfy $\left[
x^{n}\right]\mathbf{a}_{i}=0$ (since the family $\left(\mathbf{a}%
_{i}\right)_{i\in I}$ is summable). Hence, all but finitely many $i\in J$
satisfy $\left[x^{n}\right]\mathbf{a}_{i}=0$ (since $J$ is a subset of
$I$). Since we have proved this for each $n\in\mathbb{N}$, we thus conclude
that the family $\left(\mathbf{a}_{i}\right)_{i\in J}$ is summable.
\end{proof}

\begin{proposition}
\label{prop.fps.summable-sums-rule}
\lean{AlgebraicCombinatorics.FPS.summableFPS_fubini}
\leantarget
\leanok
Sums of summable families of FPSs satisfy
the usual rules for sums (such as the breaking-apart rule $\sum_{i\in S}%
a_{s}=\sum_{i\in X}a_{s}+\sum_{i\in Y}a_{s}$ when a set $S$ is the union of
two disjoint sets $X$ and $Y$). See \cite[Proposition 7.2.11]{19s} for
details. Again, the only \textbf{caveat} is about interchange
of summation signs: The equality
\[
\sum_{i\in I}\ \ \sum_{j\in J}\mathbf{a}_{i,j}=\sum_{j\in J}\ \ \sum_{i\in
I}\mathbf{a}_{i,j}
\]
holds when the family $\left(\mathbf{a}_{i,j}\right)_{\left(i,j\right)
\in I\times J}$ is summable (i.e., when for each $n\in\mathbb{N}$, all but
finitely many \textbf{pairs} $\left(i,j\right) \in I\times J$ satisfy
$\left[x^{n}\right]\mathbf{a}_{i,j}=0$); it does not generally hold if we
merely assume that the sums $\sum_{i\in I}\ \ \sum_{j\in J}\mathbf{a}_{i,j}$
and $\sum_{j\in J}\ \ \sum_{i\in I}\mathbf{a}_{i,j}$ are summable.
\end{proposition}

\begin{proof}
The proof is tedious
(as there are many rules to check), but fairly straightforward (the idea is
always to focus on a single coefficient, and then to reduce the infinite sums
to finite sums). For example, consider the discrete Fubini
rule, which says that
\[
\sum_{i\in I}\ \ \sum_{j\in J}\mathbf{a}_{i,j}=\sum_{\left(i,j\right) \in
I\times J}\mathbf{a}_{i,j}=\sum_{j\in J}\ \ \sum_{i\in I}\mathbf{a}_{i,j}
\]
whenever $\left(\mathbf{a}_{i,j}\right)_{\left(i,j\right) \in I\times
J}$ is a summable family of FPSs. In order to prove this rule, we fix a
summable family $\left(\mathbf{a}_{i,j}\right)_{\left(i,j\right) \in
I\times J}$ of FPSs. It is easy to see that the families $\left(
\mathbf{a}_{i,j}\right)_{j\in J}$ for all $i\in I$ are summable as well, as
are the families $\left(\mathbf{a}_{i,j}\right)_{i\in I}$ for all $j\in
J$, and the families $\left(\sum_{j\in J}\mathbf{a}_{i,j}\right)_{i\in I}$
and $\left(\sum_{i\in I}\mathbf{a}_{i,j}\right)_{j\in J}$.

To verify the Fubini identity, it suffices to check that
\[
\left[x^{n}\right]\left(\sum_{i\in I}\ \ \sum_{j\in J}\mathbf{a}%
_{i,j}\right) = \left[x^{n}\right]\left(\sum_{\left(i,j\right) \in
I\times J}\mathbf{a}_{i,j}\right) = \left[x^{n}\right]\left(\sum_{j\in
J}\ \ \sum_{i\in I}\mathbf{a}_{i,j}\right)
\]
for each $n\in\mathbb{N}$. Fix $n\in\mathbb{N}$; then, we have $\left[
x^{n}\right]\left(\mathbf{a}_{i,j}\right) =0$ for all but finitely many
$\left(i,j\right) \in I\times J$ (since the family $\left(\mathbf{a}%
_{i,j}\right)_{\left(i,j\right) \in I\times J}$ is summable). That is,
the set of all pairs $\left(i,j\right) \in I\times J$ satisfying $\left[
x^{n}\right]\left(\mathbf{a}_{i,j}\right) \neq0$ is finite. Hence, the
set $I^{\prime}$ of the first entries and the set $J^{\prime}$ of the second entries
of all these pairs are also finite. The three sums reduce to finite sums over
$I^{\prime}\times J^{\prime}$, which are equal by the usual Fubini rule for finite sums.
See \cite[proof of Proposition 7.2.11]{19s} for more details.
\end{proof}

\subsection{The indeterminate $x$ and powers}

\begin{definition}
\label{def.fps.x}
\lean{AlgebraicCombinatorics.FPS.X_coeff_one, AlgebraicCombinatorics.FPS.X_coeff_ne_one}
\leantarget
\leanok
Let $x$ denote the FPS $\left(0,1,0,0,0,\ldots\right)$.
In other words, let $x$ denote the FPS with $\left[x^{1}\right]x=1$ and
$\left[x^{i}\right]x=0$ for all $i\neq1$.
\end{definition}

\begin{lemma}
\label{lem.fps.xa}
\lean{AlgebraicCombinatorics.FPS.X_mul_shift, AlgebraicCombinatorics.FPS.X_mul_coeff_zero, AlgebraicCombinatorics.FPS.coeff_X_mul}
\leantarget
\leanok
Let $\mathbf{a}=\left(a_{0},a_{1},a_{2},\ldots\right)$
be an FPS. Then, $x\cdot\mathbf{a}=\left(0,a_{0},a_{1},a_{2},\ldots\right)$.
\end{lemma}

\begin{proof}
\leanok
If $n$ is a positive integer, then
\begin{align*}
\left[x^{n}\right]\left(x\cdot\mathbf{a}\right) & =\sum_{i=0}%
^{n}\left[x^{i}\right]x\cdot a_{n-i}\\
& =\underbrace{\left[x^{0}\right]x}_{=0}\cdot\,a_{n-0}+\underbrace{\left[x^{1}\right]x}%
_{=1}\cdot\,a_{n-1}+\sum_{i=2}^{n}\underbrace{\left[x^{i}\right]
x}_{\substack{=0\\\text{(since }i\geq2>1\text{)}}}\cdot\,a_{n-i}\\
& =a_{n-1}.
\end{align*}
A similar argument for $n=0$ yields $\left[x^{0}\right]\left(x\cdot\mathbf{a}\right)=0$.
Thus, for each $n\in\mathbb{N}$,
\[
\left[x^{n}\right]\left(x\cdot\mathbf{a}\right) =
\begin{cases}
a_{n-1}, & \text{if }n>0;\\
0, & \text{if }n=0.
\end{cases}
\]
In other words, $x\cdot\mathbf{a}=\left(0,a_{0},a_{1},a_{2},\ldots\right)$.
\end{proof}

\begin{proposition}
\label{prop.fps.xk}
\lean{AlgebraicCombinatorics.FPS.X_pow_coeff}
\leantarget
\leanok
We have
\[
x^{k}=\left(\underbrace{0,0,\ldots,0}_{k\text{ times}},1,0,0,0,\ldots
\right) \ \ \ \ \ \ \ \ \ \ \text{for each }k\in\mathbb{N}.
\]
\end{proposition}

\begin{proof}
We induct on $k$.

\textit{Induction base:} We have $x^{0}=\underline{1}=\left(1,0,0,0,0,\ldots
\right) = \left(\underbrace{0,0,\ldots,0}_{0\text{ times}},1,0,0,0,\ldots
\right)$.

\textit{Induction step:} Let $m\in\mathbb{N}$. Assume that Proposition
\ref{prop.fps.xk} holds for $k=m$. We have $x^{m}=\left(\underbrace{0,0,\ldots,0}_{m\text{ times}%
},1,0,0,0,\ldots\right)$. Thus, Lemma \ref{lem.fps.xa} (applied to $\mathbf{a}=x^{m}$) yields
\[
x\cdot x^{m}=\left(0,\underbrace{0,0,\ldots,0}_{m\text{ times}%
},1,0,0,0,\ldots\right) = \left(\underbrace{0,0,\ldots,0}_{m+1\text{ times}%
},1,0,0,0,\ldots\right).
\]
Since $x\cdot x^{m}=x^{m+1}$, this completes the induction step.
\end{proof}

\begin{corollary}
\label{cor.fps.sumakxk}
\lean{AlgebraicCombinatorics.FPS.fps_eq_tsum_coeff, AlgebraicCombinatorics.FPS.summableFPS_monomial_family}
\leantarget
\leanok
Any FPS $\left(a_{0},a_{1},a_{2},\ldots\right) \in
K\left[\left[x\right]\right]$ satisfies
\[
\left(a_{0},a_{1},a_{2},\ldots\right) = a_{0}+a_{1}x+a_{2}x^{2}+a_{3}%
x^{3}+\cdots = \sum_{n\in\mathbb{N}}a_{n}x^{n}.
\]
In particular, the right hand side here is well-defined, i.e., the family
$\left(a_{n}x^{n}\right)_{n\in\mathbb{N}}$ is summable.
\end{corollary}

\begin{proof}
By Proposition \ref{prop.fps.xk}, we have
\begin{align*}
& a_{0}+a_{1}x+a_{2}x^{2}+a_{3}x^{3}+\cdots\\
& =\ \ \ \ a_{0}\left(1,0,0,0,\ldots\right) \\
& \ \ \ \ +a_{1}\left(0,1,0,0,\ldots\right) \\
& \ \ \ \ +a_{2}\left(0,0,1,0,\ldots\right) \\
& \ \ \ \ +a_{3}\left(0,0,0,1,\ldots\right) \\
& \ \ \ \ +\cdots\\
& =\ \ \ \ \left(a_{0},0,0,0,\ldots\right) \\
& \ \ \ \ +\left(0,a_{1},0,0,\ldots\right) \\
& \ \ \ \ +\left(0,0,a_{2},0,\ldots\right) \\
& \ \ \ \ +\left(0,0,0,a_{3},\ldots\right) \\
& \ \ \ \ +\cdots\\
& =\left(a_{0},a_{1},a_{2},a_{3},\ldots\right).
\end{align*}
\end{proof}

\subsection{Helper lemmas from the formalization}

\begin{lemma}
\label{lem.fps.ext}
\lean{AlgebraicCombinatorics.FPS.fps_ext_iff}
\leanhelper
\leanok
Two FPSs $f$ and $g$ are equal if and only if $\left[x^{n}\right]f = \left[x^{n}\right]g$
for all $n\in\mathbb{N}$.
\end{lemma}

\begin{proof}
\leanok
Immediate from the definition of FPS as a sequence.
\end{proof}

\begin{lemma}
\label{lem.fps.coeff.zero.mul}
\lean{AlgebraicCombinatorics.FPS.coeff_zero_mul_fps}
\leanhelper
\leanok
For any $\mathbf{a},\mathbf{b}\in K\left[\left[x\right]\right]$, we have
$\left[x^{0}\right]\left(\mathbf{ab}\right) = \left[x^{0}\right]
\mathbf{a}\cdot\left[x^{0}\right]\mathbf{b}$.
That is, the constant term of a product is the product of the constant terms.
\end{lemma}

\begin{proof}
\leanok
Follows from the product formula with $n=0$: $\left[x^{0}\right](\mathbf{ab}) = \sum_{i=0}^{0} \left[x^{i}\right]\mathbf{a}\cdot\left[x^{0-i}\right]\mathbf{b} = \left[x^{0}\right]\mathbf{a}\cdot\left[x^{0}\right]\mathbf{b}$.
\end{proof}

\begin{lemma}
\label{lem.fps.xk.mul.shift}
\lean{AlgebraicCombinatorics.FPS.coeff_X_pow_mul}
\leanhelper
\leanok
For any $k\in\mathbb{N}$ and any FPS $\mathbf{a}$,
\[
\left[x^{n}\right]\left(x^{k}\cdot\mathbf{a}\right) =
\begin{cases}
0, & \text{if }n<k;\\
\left[x^{n-k}\right]\mathbf{a}, & \text{if }n\geq k.
\end{cases}
\]
That is, multiplication by $x^{k}$ shifts coefficients by $k$ positions.
\end{lemma}

\begin{proof}
\leanok
By induction on $k$, using Lemma \ref{lem.fps.xa} for the induction step.
\end{proof}

\begin{lemma}
\label{lem.fps.summable.add}
\lean{AlgebraicCombinatorics.FPS.summableFPS_add}
\leanhelper
\leanok
If $\left(\mathbf{a}_{i}\right)_{i\in I}$ and $\left(\mathbf{b}_{i}\right)_{i\in I}$
are summable families of FPSs, then so is $\left(\mathbf{a}_{i}+\mathbf{b}_{i}\right)_{i\in I}$.
\end{lemma}

\begin{proof}
\leanok
For each $n\in\mathbb{N}$,
$\{i \mid \left[x^{n}\right](\mathbf{a}_{i}+\mathbf{b}_{i})\neq 0\}
\subseteq \{i \mid \left[x^{n}\right]\mathbf{a}_{i}\neq 0\} \cup
\{i \mid \left[x^{n}\right]\mathbf{b}_{i}\neq 0\}$, and the union of two
finite sets is finite.
\end{proof}

\begin{lemma}
\label{lem.fps.summable.neg}
\lean{AlgebraicCombinatorics.FPS.summableFPS_neg}
\leanhelper
\leanok
If $\left(\mathbf{a}_{i}\right)_{i\in I}$ is a summable family of FPSs,
then so is $\left(-\mathbf{a}_{i}\right)_{i\in I}$.
\end{lemma}

\begin{proof}
\leanok
For each $n$, $\left[x^{n}\right](-\mathbf{a}_{i})=-\left[x^{n}\right]\mathbf{a}_{i}$,
so $\{i \mid \left[x^{n}\right](-\mathbf{a}_{i})\neq 0\} = \{i \mid \left[x^{n}\right]\mathbf{a}_{i}\neq 0\}$.
\end{proof}

\begin{lemma}
\label{lem.fps.essfin.add}
\lean{AlgebraicCombinatorics.FPS.essentiallyFinite_add}
\leanhelper
\leanok
The sum of two essentially finite families is essentially finite.
\end{lemma}

\begin{proof}
\leanok
$\{i \mid a_{i}+b_{i}\neq 0\}\subseteq \{i \mid a_{i}\neq 0\}\cup\{i \mid b_{i}\neq 0\}$,
and the union of two finite sets is finite.
\end{proof}

\begin{lemma}
\label{lem.fps.essfin.neg}
\lean{AlgebraicCombinatorics.FPS.essentiallyFinite_neg}
\leanhelper
\leanok
The negation of an essentially finite family is essentially finite.
\end{lemma}

\begin{proof}
\leanok
$\{i \mid -a_{i}\neq 0\}=\{i \mid a_{i}\neq 0\}$.
\end{proof}

\begin{lemma}
\label{lem.fps.essfin.sub}
\lean{AlgebraicCombinatorics.FPS.essentiallyFinite_sub}
\leanhelper
\leanok
The difference of two essentially finite families is essentially finite.
\end{lemma}

\begin{proof}
\leanok
Write $a_i - b_i = a_i + (-b_i)$ and apply
Lemma~\ref{lem.fps.essfin.add} and Lemma~\ref{lem.fps.essfin.neg}.
\end{proof}

\begin{definition}
\label{def.fps.essFinSum}
\lean{AlgebraicCombinatorics.FPS.essFinSum}
\leanhelper
\leanok
For an essentially finite family $(a_i)_{i\in I}$, the \emph{essentially finite sum}
$\sum_{i\in I} a_i$ is defined as $\sum_{i\in S} a_i$ where $S = \{i\in I \mid a_i \neq 0\}$.
\end{definition}

\begin{lemma}
\label{lem.fps.essFinSum.add}
\lean{AlgebraicCombinatorics.FPS.essFinSum_add}
\leanhelper
\leanok
The essentially finite sum distributes over addition:
if $(a_i)_{i\in I}$ and $(b_i)_{i\in I}$ are essentially finite families,
then $\sum_{i\in I}(a_i + b_i) = \sum_{i\in I} a_i + \sum_{i\in I} b_i$.
\end{lemma}

\begin{proof}
\leanok
All three sums can be computed over the finite set
$S = \{i \mid a_i\neq 0\}\cup\{i \mid b_i\neq 0\}$,
so this reduces to the finite sum identity $\sum_{i\in S}(a_i+b_i)=\sum_{i\in S}a_i+\sum_{i\in S}b_i$.
\end{proof}

\begin{lemma}
\label{lem.fps.essFinSum.smul}
\lean{AlgebraicCombinatorics.FPS.essFinSum_smul}
\leanhelper
\leanok
For a scalar $c\in K$ and an essentially finite family $(a_i)_{i\in I}$,
$\sum_{i\in I} c\, a_i = c\cdot\sum_{i\in I} a_i$.
\end{lemma}

\begin{proof}
\leanok
The support of $(c\,a_i)$ is contained in the support of $(a_i)$,
so this reduces to the finite identity $\sum_{i\in S}c\,a_i = c\sum_{i\in S}a_i$.
\end{proof}

\begin{lemma}
\label{lem.fps.summable.sub}
\lean{AlgebraicCombinatorics.FPS.summableFPS_sub}
\leanhelper
\leanok
If $\left(\mathbf{a}_{i}\right)_{i\in I}$ and $\left(\mathbf{b}_{i}\right)_{i\in I}$
are summable families of FPSs, then so is $\left(\mathbf{a}_{i}-\mathbf{b}_{i}\right)_{i\in I}$.
\end{lemma}

\begin{proof}
\leanok
Write $\mathbf{a}_i - \mathbf{b}_i = \mathbf{a}_i + (-\mathbf{b}_i)$
and apply Lemma~\ref{lem.fps.summable.add} and Lemma~\ref{lem.fps.summable.neg}.
\end{proof}

\begin{lemma}
\label{lem.fps.summable.of_essfin}
\lean{AlgebraicCombinatorics.FPS.summableFPS_of_essentiallyFinite}
\leanhelper
\leanok
If $(\mathbf{a}_i)_{i\in I}$ is a family of FPSs such that
$\{i \mid \mathbf{a}_i \neq 0\}$ is finite, then the family is summable.
\end{lemma}

\begin{proof}
\leanok
For each $n$, $\{i \mid [x^n]\mathbf{a}_i \neq 0\} \subseteq \{i \mid \mathbf{a}_i \neq 0\}$,
which is finite by hypothesis.
\end{proof}

\begin{lemma}
\label{lem.fps.coeff_monomial_X_pow}
\lean{AlgebraicCombinatorics.FPS.coeff_monomial_mul_X_pow}
\leanhelper
\leanok
For $a\in K$ and $k,n\in\mathbb{N}$,
$[x^n](\underline{a}\cdot x^k) = \begin{cases} a & \text{if } n = k, \\ 0 & \text{otherwise.}\end{cases}$
\end{lemma}

\begin{proof}
\leanok
Follows from $[x^n]x^k = \delta_{n,k}$ (Proposition~\ref{prop.fps.xk})
and $[x^n](\underline{a}\cdot f) = a\cdot[x^n]f$.
\end{proof}

\begin{lemma}
\label{lem.fps.mul_X_shift}
\lean{AlgebraicCombinatorics.FPS.mul_X_shift}
\leanhelper
\leanok
Multiplying on the right by $x$ also shifts coefficients:
$[x^{n+1}](\mathbf{a}\cdot x) = [x^n]\mathbf{a}$.
\end{lemma}

\begin{proof}
\leanok
By commutativity of multiplication, $\mathbf{a}\cdot x = x\cdot\mathbf{a}$,
so this follows from Lemma~\ref{lem.fps.xa}.
\end{proof}

\begin{lemma}
\label{lem.fps.mul_X_pow_shift}
\lean{AlgebraicCombinatorics.FPS.coeff_mul_X_pow}
\leanhelper
\leanok
For any $k\in\mathbb{N}$ and any FPS $\mathbf{a}$,
$[x^n](\mathbf{a}\cdot x^k) = \begin{cases} 0 & \text{if }n<k, \\ [x^{n-k}]\mathbf{a} & \text{if }n\geq k.\end{cases}$
\end{lemma}

\begin{proof}
\leanok
By commutativity: $\mathbf{a}\cdot x^k = x^k\cdot\mathbf{a}$,
so this follows from Lemma~\ref{lem.fps.xk.mul.shift}.
\end{proof}

\subsection{Generating functions}

\begin{definition}
\label{def.fps.gf}
\lean{AlgebraicCombinatorics.FPS.generatingFunction}
\leanhelper
\leanok
The \emph{(ordinary) generating function} of a sequence $(a_0, a_1, a_2, \ldots)$
is the FPS $(a_0, a_1, a_2, \ldots) = \sum_{n\geq 0} a_n x^n$.
\end{definition}

\begin{lemma}
\label{lem.fps.gf.coeff}
\lean{AlgebraicCombinatorics.FPS.generatingFunction_coeff}
\leanhelper
\leanok
The $n$-th coefficient of the generating function of a sequence $(a_n)$ is $a_n$.
\end{lemma}

\begin{proof}
\leanok
Immediate from the definitions.
\end{proof}

\subsection{The Chu--Vandermonde identity}

\begin{proposition}
\label{prop.binom.vandermonde.NN}
\lean{AlgebraicCombinatorics.FPS.vandermonde_nat}
\leantarget
\leanok
Let $a,b\in\mathbb{N}$, and let
$n\in\mathbb{N}$. Then,
\[
\binom{a+b}{n}=\sum_{k=0}^{n}\binom{a}{k}\binom{b}{n-k}.
\]
\end{proposition}

\begin{proof}
Multiply out the identity $\left(1+x\right)^{a+b}=\left(1+x\right)^{a}\left(1+x\right)^{b}$
using the binomial formula and compare the coefficient of $x^{n}$ on both sides.
The left side gives $\binom{a+b}{n}$, and the right side gives
$\sum_{k=0}^{n}\binom{a}{k}\binom{b}{n-k}$ by the product formula for FPSs.
\end{proof}

\begin{lemma}
\label{lem.fps.vandermonde_nat_range}
\lean{AlgebraicCombinatorics.FPS.vandermonde_nat'}
\leanhelper
\leanok
The Vandermonde identity in range-sum form:
for $a, b, n \in \mathbb{N}$,
$\dbinom{a+b}{n} = \sum_{k=0}^{n}\dbinom{a}{k}\dbinom{b}{n-k}$.
\end{lemma}

\begin{proof}
\leanok
Rewrites the antidiagonal sum from Proposition~\ref{prop.binom.vandermonde.NN}
as a range sum.
\end{proof}

\begin{theorem}[Vandermonde convolution identity, aka Chu--Vandermonde identity]
\label{thm.binom.vandermonde.CC}
\lean{AlgebraicCombinatorics.FPS.chuVandermonde}
\leantarget
\leanok
Let $a,b\in\mathbb{C}$, and let
$n\in\mathbb{N}$. Then,
\[
\binom{a+b}{n}=\sum_{k=0}^{n}\binom{a}{k}\binom{b}{n-k}.
\]
\end{theorem}

\begin{proof}
Fix $n\in\mathbb{N}$ and $b\in\mathbb{N}$. By Proposition
\ref{prop.binom.vandermonde.NN}, the equality
\[
\binom{a+b}{n}=\sum_{k=0}^{n}\binom{a}{k}\binom{b}{n-k}
\]
holds for each $a\in\mathbb{N}$. However, both sides are
polynomials in $a$:
\begin{align*}
\binom{a+b}{n} & =\frac{\left(a+b\right)\left(a+b-1\right)\cdots\left(a+b-n+1\right)}{n!},\\
\sum_{k=0}^{n}\binom{a}{k}\binom{b}{n-k} & =\sum_{k=0}^{n}\frac{a\left(
a-1\right)\cdots\left(a-k+1\right)}{k!}\binom{b}{n-k}.
\end{align*}
Two univariate polynomials that agree on all of $\mathbb{N}$ (infinitely many points)
must be identical. Hence the equality holds for all $a\in\mathbb{C}$.

Now fix $a\in\mathbb{C}$ and repeat the argument with $b$ as the variable:
both sides are polynomials in $b$ that agree on $\mathbb{N}$, hence they agree
on $\mathbb{C}$.
(See \cite[proofs of Lemma 7.5.8 and Theorem 7.5.3]{20f} or
\cite[\S 2.17.3]{19s} for this proof in more detail. Alternatively, see
\cite[\S 3.3.2 and \S 3.3.3]{detnotes} for two other proofs.)
\end{proof}
